\chapter{Draumen om ein designvitskap}

\begin{motto}
Ein mann utanfrå gav faget den teorien det hadde bede om i hundre år.\\ Faget tok imot prestisjen og lét teorien liggje.
\end{motto}

Det finst eitt augneblink i designhistoria då faget vart tilbode nett det det heile tida hadde sagt at det mangla: ein rasjonell, formell kjerne, ein måte å tenkje om formgjeving på som kunne stå ved sida av matematikken og fysikken utan å raudne. Tilbodet kom ikkje innanfrå. Det kom frå ein økonom og kognisjonsvitar ved Carnegie Mellon, ein mann som aldri teikna ein stol og aldri sette sin fot i eit atelier som student, og som likevel skreiv den mest ambisiøse setninga nokon har skrive om kva design kunne vere. Herbert Simon kalla design ein \emph{science of the artificial} — ein vitskap om det kunstige, om det som ikkje er gjeve av naturen, men som kunne ha vore annleis \autocite{simon1969sciences}. I den setninga låg eit heilt program. Og historia om kvifor programmet ikkje feste seg i faget det var meint å redde, er ei av dei reinaste illustrasjonane denne boka har på den dobbelte bokføringa: faget tok imot statusen Simon gav det, og avviste disiplinen som var prisen for statusen.

Boka kjem ikkje til Simon tilfeldig på dette punktet. Det førre kapitlet følgde metoderørsla — det optimistiske forsøket frå femti- og sekstitalet på å gjere design til ei systematisk, stegvis, mest mogleg rasjonell verksemd, med diagram, fasar og prosedyrar. Den rørsla var i emning over heile den vestlege designverda, frå London til Ulm, og ho var driven av ei sann von om at formgjeving kunne bli noko meir disiplinert enn smak. Men ho mangla eit teoretisk fundament; ho hadde metodar utan ein teori om kvifor metodane skulle verke, prosedyrar utan ein modell av kva ein prosedyre eigentleg gjorde med eit designproblem. Simon er den som kjem og tilbyr nett den manglande undergrunnen. Der metoderørsla hadde framgangsmåtar, gav Simon ein teori om kva ein framgangsmåte er: ein måte å søkje gjennom eit rom av moglege artefakt på. Han er ikkje ein avstikkar frå metoderørsla; han er dens mogelege fullbyrding, den teoretiske ryggrada ho aldri sjølv greidde å reise. Det gjer det desto meir talande at faget tok imot rørsla og ikkje ryggrada.

\section{Det kunstige som eige kunnskapsobjekt}

Ein må forstå kor radikalt og kor frigjerande Simons grep var, for å forstå kor stort sviket mot det var. Naturvitskapen studerer det som er: det som finst uavhengig av menneskeleg vilje, og som ikkje kunne ha vore annleis enn lovene gjer det. Men kring oss ligg ei heil verd som ikkje er slik. Stolen, byen, datamaskinen, organisasjonen, læreplanen — alt dette er kunstig i ei presis tyding: det er menneskeskapt, det tener eit føremål, og det kunne ha vore utforma annleis. Simons innsikt var at dette kunstige ikkje berre er eit rot av tilfeldige menneskeval som ligg under terskelen for vitskap. Det har sin eigen logikk, og den logikken kan studerast strengt. Vitskapen om det kunstige er vitskapen om korleis indre struktur vert tilpassa ytre føremål gjennom eit grensesnitt — og design, sa Simon, er nett denne tilpassinga. «Alle som tenkjer ut handlemåtar som endrar eksisterande situasjonar til føretrekte, driv med design,» skreiv han, og med den setninga løfta han design frå eit yrke til ein allmenn intellektuell operasjon, like grunnleggjande som analyse.

Dette var ei gåve faget hadde drøymt om utan å kunne formulere. I to kapittel har vi sett faget streve med ein tom formel og ein taus kompetanse; her kom ein som tilbaud det objektet ein teori treng. For Simons design var ikkje ein stil og ikkje eit handlag. Det var ein klasse av problem med ein delt struktur, og difor noko ein kunne ha ein generell teori om. Han skilde med vilje mellom naturvitskapens spørsmål — korleis er verda? — og designvitskapens spørsmål — korleis bør artefaktet vere, gjeve eit mål? Det fyrste handlar om det nødvendige, det andre om det moglege. Og han hevda at det moglege kunne studerast like systematisk som det nødvendige, så lenge ein var villig til å byggje apparatet for det. Apparatet vert som regel kalla det fyrste alvorlege framlegget om eit fundament for design som verkeleg fortente namnet, og det er verdt å sjå kva det inneheldt.

\section{Avgjerd, søk, tilfredsstilling}

Tre omgrep ber Simons konstruksjon, og dei er verdt å namne kvar for seg, fordi det er presis desse faget ikkje tok opp. Det fyrste er \emph{avgjerd}. Design er for Simon i kjernen ei rekkje val under avgrensingar: gjeve eit mål og eit sett vilkår, kva handling fører nærast målet? Det andre er \emph{søk}. Mengda av moglege artefakt er enorm, oftast altfor stor til å granskast uttømmande, og difor er design ikkje ei rekning men ei \emph{leiting} gjennom eit rom av alternativ, styrt av heuristikkar som gjer leitinga handterleg. Forma er ein posisjon i eit moglegheitsrom; designprosessen er ein veg gjennom det rommet. Det tredje, og det mest karakteristiske for Simon, er \emph{tilfredsstilling} — hans \textit{satisficing}. Sidan det fullkomne søket er umogleg, leitar den verkelege formgjevaren ikkje etter den optimale løysinga, men etter ei løysing som er god nok mot ein terskel. Rasjonaliteten er avgrensa, ikkje fordi mennesket er dumt, men fordi problema er for store og tida for knapp; og ein teori om design må vere ein teori om denne avgrensa, søkjande, tilfredsstillande rasjonaliteten.

Ein kunne innvende at design sjeldan ser slik ut innanfrå — at formgjevaren ikkje opplever seg som ein søkjealgoritme. Simon såg innvendinga komme og svara på henne i den teksten som meir enn nokon annan strekkjer programmet hans mot røynda av verkeleg formgjeving. I \textit{The Structure of Ill-Structured Problems} tek han fatt i nett den eigenskapen ved designproblem som faget seinare skulle bruke som unnskyldning for å sleppe teori: at dei er dårleg strukturerte, utan klar formulering, utan eit lukka rom av alternativ gjeve på førehand \autocite{simon1973illstructured}. Og her er det avgjerande: Simon brukar ikkje dette som bevis på at design unndreg seg formell handsaming. Han brukar det som spesifikasjonen på kva ein teori om design må makte. Hans påstand er at eit dårleg strukturert problem vert løyst ved at det vert gjort velstrukturert i bitar, lokalt og mellombels, gjennom den same søkjeprosessen — at strukturen ikkje er gjeven, men vert frambrakt undervegs. Med andre ord: at problemet er vrangt, er ikkje slutten på teorien; det er det teorien skal forklare. Hugs denne setninga. Vi skal sjå faget velje det stikk motsette.

Det er verdt å dvele eit augneblink ved kor sterk denne posisjonen eigentleg er, fordi faget seinare karikerte henne for å lettare avvise henne. Karikaturen seier at Simon trudde design var ei rein optimering, ei rekneoppgåve med eitt rett svar som berre venta på å bli rekna ut. Det er det motsette av kva han skreiv. Heile poenget med satisficing er at det \emph{ikkje} finst noko utreknbart optimum innan rekkjevidd, at formgjevaren arbeider i blinde gjennom eit rom han ikkje kan oversjå, og at dette er normaltilstanden, ikkje eit unntak. Heile poenget med det dårleg strukturerte problemet er at problemet ikkje kjem ferdig formulert, at sjølve formuleringa er ein del av arbeidet. Simon var langt nærare den verkelege erfaringa av å designe enn motstandarane hans ville vedgå; han mangla ikkje sansen for at formgjeving er rotete, usikker og famlande. Han hadde berre, til skilnad frå faget, den ambisjonen at det rotete, usikre og famlande sjølv kunne bli gjenstand for ein presis teori. Det er denne ambisjonen, ikkje ein påstått kald rasjonalisme, faget eigentleg sa nei til.

\section{Prestisjen utan disiplinen}

No til det avgjerande spørsmålet, og det ubehagelege svaret. Simon gav faget eit objekt, eit apparat og ein ambisjon på vitskapleg nivå. Kva gjorde faget med det? Det enkle svaret er: lite. Det meir presise svaret er at Simons program slo rot, men ikkje i designskulane. Det vart adoptert, vidareutvikla og gjort fruktbart av informatikken, der søk og heuristikk vart sjølve grunnvollen for kunstig intelligens; av økonomien, der avgrensa rasjonalitet og satisficing omforma teorien om val; og av kognisjonsvitskapen, der mennesket som informasjonshandsamande problemløysar vart eit heilt forskingsprogram. Simon fekk Nobelprisen i økonomi og Turingprisen i informatikk. Han fekk ingen pris frå designfaget, og — meir talande — designfaget bygde ikkje vidare på han. Apparatet hans vart kjernen i tre andre disiplinar og ein fotnote i sin eigen.

Kvifor? Den vanlege forklaringa innanfor faget er at Simon var for snever — at hans rasjonalisme ikkje fanga det intuitive, kroppslege og verdiladde ved verkeleg formgjeving. Det er ikkje heilt usant, og vi kjem til kritikken straks. Men som forklaring på fråvêret er det utilstrekkeleg, for det forklarar berre kvifor faget kunne hatt innvendingar, ikkje kvifor det ikkje bygde noko alternativt på same nivå. Den ærlege forklaringa er hardare. Å adoptere Simons disiplin ville kravd at faget gjorde noko det aldri har vore villig til: å formalisere, å gjere påstandane sine etterprøvbare, å byggje kumulativt, å la seg korrigere. Det ville kravt matematikk der faget ville ha atelier, og felles omgrep der faget ville ha personleg dømekraft. Simons gåve kom med ein rekning, og rekninga var disiplin. Faget tok imot \emph{prestisjen} — gleda over å kunne seie at design er ein vitskap om det kunstige, at formgjeving er ein rasjonell operasjon på linje med analyse — og let \emph{disiplinen} liggje att i pakken. Dette er den same operasjonen vi alt har sett: faget vil ha vitskapens stilling utan vitskapens forpliktingar. Simon gjorde berre prisen for éin gong eksplisitt, og faget såg prisen og takka nei, samstundes som det heldt på kvitteringa for å vise fram.

Du ser konsekvensen i korleis Simon vert undervist, der han i det heile vert undervist. Han vert sitert for setninga om å endre situasjonar til føretrekte — den prestisjeberande, allmenne, smigrande setninga — og sjeldan for apparatet som gjorde setninga til meir enn ein aforisme. Søket, heuristikkane, satisficing, den formelle handsaminga av det dårleg strukturerte: det som var disiplinen, fell bort. Att står slagordet. Faget gjorde med Simon det same som det gjorde med Sullivan: det reduserte ein potensiell grunnvoll til eit motto det kunne ha på veggen.

Og legg merke til kva faget bygde i staden, for det er her fråvêret tek positiv form. I staden for å ta opp Simons formelle program, formulerte faget på åttitalet ein motposisjon som vart langt meir innflytelsesrik innanfor designskulane enn Simon nokon gong vart: tanken om at design har sine \emph{eigne} måtar å vite på, ein eigen kunnskapsform som ikkje let seg redusere til vitskapens. Nigel Cross gav denne tanken sitt kanoniske uttrykk då han hevda at det finst ein eigen designarleg måte å kjenne verda på, jamstelt med, men ulik, dei to kulturane i vitskap og humaniora \autocite{cross1982}. Eg vil vere rettferdig mot Cross, for han var seriøs og las Simon nøye, og det er sant at formgjeving har eigenarta trekk. Men sjå kva grepet gjer i fagets sjølvforståing. Der Simon sa: design er ein vitskap, og her er apparatet — kjem motposisjonen og seier: design er sin eigen kunnskapsform, med eigne kriterium. Det høyrest ut som ei sjølvstendiggjering. I praksis vart det ei orsaking. For «eigne måtar å vite på» kom aldri med eit apparat som Simons; det kom ikkje med etterprøvbare påstandar, kumulative resultat eller felles omgrep. Det kom med ein påstand om at faget ikkje treng slikt, fordi det veit på sin eigen måte. Faget valde altså ikkje mellom Simons vitskap og inga grunngjeving; det valde ei tredje form, som hadde vitskapens sjølvtillit og handverkets fråvær av forpliktingar — ei kunnskapsteoretisk grunngjeving for å sleppe å grunngje seg. Mellom det formelle programmet og det designarlege vitet valde faget det designarlege, og kalla valet ei frigjering, når det var ei flukt.

\section{Kuhn, og kvifor poda ikkje kunne festne}

For å sjå at dette ikkje berre var vrangvilje eller doven, men noko strukturelt, lønner det seg å låne eit blikk frå vitskapshistoria sjølv. Same tiåret som Simon skreiv, gav Thomas Kuhn faget eit omgrep det kunne ha brukt til å forstå si eiga stode, om det hadde våga: skiljet mellom det \emph{før-paradigmatiske} og det \emph{paradigmatiske} \autocite{kuhn1962}. Ein moden vitskap, hjå Kuhn, kviler på eit paradigme: eit delt sett av eksemplariske problemløysingar, omgrep og metodar som heile fellesskapet tek for gjeve, og som let forskarane byggje på kvarandre i staden for å byrje på nytt kvar gong. Eit før-paradigmatisk felt manglar dette. Der finst skular i staden for ein disiplin, kvar med sine eigne grunnomgrep, sine eigne autoritetar, sine eigne kriterium for kva som tel som godt arbeid — og difor inga sams grunn å samle innsats på.

Sjå designfaget gjennom det blikket, og biletet vert smerteleg klart. Faget er, i Kuhns tyding, før-paradigmatisk, og har vore det heile tida. Det har ingen delt kjerne av eksemplariske løysingar alle reknar med; det har atelier som er usamde om sjølve spørsmålet om design har ein metode. Og her ligg grunnen til at Simons pode ikkje kunne festne: eit paradigme kan ikkje takast imot av eit felt som ikkje har den institusjonelle og intellektuelle infrastrukturen til å bere eit paradigme i det heile. For å adoptere ein vitskapleg kjerne må eit fag alt vere organisert som ein vitskap — det må ha tidsskrift som krev etterprøvbarheit, ein utdanning som lærer ein delt kanon før personleg uttrykk, ein praksis for å forkaste det som er motbevist. Designfaget hadde ingen av delane. Det kunne difor ikkje ta imot Simon sjølv om det hadde villa, like lite som ein kropp utan immunsystem kan ta imot ein vaksine. Kuhn lèt oss seie det presist: Simon tilbaud eit paradigme til eit felt som var konstitusjonelt ute av stand til å hyse eitt. Prestisjen kunne festne, fordi prestisje ikkje krev infrastruktur. Disiplinen kunne ikkje, fordi disiplin er infrastruktur.

Dette gjev oss òg den rette plassen for Simon i historia vår. Han er ikkje den tenkjaren faget følgde og gjekk vill med; han er den tenkjaren faget \emph{ikkje} følgde, og som difor måler avstanden mellom det faget kunne ha vore og det det vart. Bruce Archer og dei som ville gjere design til ein disiplin på sekstitalet, hadde lese Simon og drøymt hans draum; arven frå Ulm peika same veg, mot ein formgjeving grunngjeven på metode og kunnskap heller enn på smak. Men der Ulm fall på indre strid og ytre press, fall Simons program på noko djupare: på at mottakaren ikkje var bygd for gåva. Lina frå Ulm gjennom Archer til Simon er lina der faget kom nærast eit fundament — og difor lina der fråvêret av fundament er lettast å lese, fordi vi her kan sjå presis kva som vart tilbode og presis kva som vart lagt bort.

\vignett

Simon hadde sagt at det vrange problemet var det teorien skulle forklare. Det neste store grepet i faget tek same kjensgjerning — at designproblem er vrange — og dreg den motsette slutninga: at vrangskapen viser at teori er umogleg, og at krava til vitskap difor er malplasserte. Reaksjonen mot Simons rasjonalisme har eit namn, og det er der vi går no.

\vidarelesing{\fullcite{simon1969sciences}; \fullcite{simon1973illstructured}; \fullcite{kuhn1962}; \fullcite{archer1979discipline}; \fullcite{cross1982}; \fullcite{rittel1973}.}
